\section{本章小结}
CPU 主机负责组织任务，GPU 设备负责并行计算。驱动负责访问设备，Toolkit 提供编译工具和开发库；版本信息要分开检查。应熟练计算网格索引、检查越界、核对每次 API 返回值，并计入拷贝与同步开销。共享内存和推理引擎属于后续优化入口，先保证输出正确。

\begin{problemset}
\item 基础：阈值化每个像素与按前一元素结果递推，哪一个更容易直接映射为独立线程？
\item 基础：nvidia-smi 正常但 nvcc 不存在，能否断定驱动未安装？
\item 基础：两个长度 512 的 float 输入和一个输出，一次流程需拷贝多少字节？
\item 基础：n=513、每块 256 线程，需要多少块？最后有多少线程必须跳过访问？
\item 基础：Shared memory 的共享范围是什么？部分线程在屏障之前返回有什么风险？
\item 提高：kernel 已成功提交，程序是否就一定计算成功？如何捕获执行错误？
\item 实践：有可用 GPU 时分别记录 nvidia-smi 与 nvcc 的版本；无 GPU 时分析教师提供的这两段输出，说明各自代表什么。
\end{problemset}

\section{本章验收}
\begin{enumerate}
    \item 展示现有驱动与 Toolkit 状态，解释安装向导中 Driver 选项的处理依据。
    \item 用 nvcc 或 CMake 编译向量加法，对检查者给出的非整块长度验证全部结果，并展示错误检查位置。
    \item 测量包含数据传输的完整流程耗时；选做在固定块大小下改写共享内存示例并验证边界长度。
\end{enumerate}
